\documentclass[10pt,a4paper]{article}
\usepackage{fontspec,amsmath,amssymb,graphicx}
\usepackage[margin=18mm]{geometry}
\setmainfont{Times New Roman}
\setlength{\parindent}{0pt}
\setlength{\parskip}{5pt}
\allowdisplaybreaks
\sloppy
\emergencystretch=2em
\begin{document}
\section*{BT\_5\_4\_3\_10 — Lời giải bổ sung}
Viết ma trận $T=A$. Với mọi ma trận đối xứng có vết 0, đa thức $u(y)=y^THy$ điều hòa. Giả thiết cho $0=\Delta(u\circ A)=2\operatorname{tr}(HAA^T)$. Trực giao với mọi $H$ vết 0 buộc $AA^T=c^2I$, $c\ge0$. Nếu $c=0$ thì $A=0$; nếu $c>0$, $A=cQ$ với $Q$ trực giao. Đây là kết luận, kể cả phép biến đổi 0.\end{document}